\clearpage
\fancyhead[L]{\small PALLAS.}
\hypertarget{native-267-lower}{}
\begin{center}
PALLAS.

\rule{0.5\linewidth}{0.4pt}

Monatliche Correspondenz. Band XI. S. 376–383. 1805 April.

\rule{0.5\linewidth}{0.4pt}
\end{center}

Zwar konnte \emph{Pallas} schon seit mehrern Monaten nicht mehr beobachtet werden; allein der Fleiss
des Dr. \textsc{Gauss} setzt uns in Stand, unsern Lesern hier noch einiges über diesen neuen Planeten in theoretischer
Hinsicht mittheilen zu können, was vorzüglich dazu dienen wird, dessen Wiederauffindung bei
der nächsten, wahrscheinlich etwas entfernten Epoche von Sichtbarkeit zu erleichtern. Dr. \textsc{Gauss}, dessen
unermüdeter Arbeitsamkeit und anhaltendem Streben nach Vervollkommnung der berechneten Planetenbahnen
nur Mangel an Stoff Grenzen zu setzen vermag, äusserte mehrmal seine Unzufriedenheit
über die in so geringer Anzahl vorhandenen Beobachtungen der Pallas, die es ihm unmöglich machte,
eine fernere Berichtigung der VII. Elemente dieses Planeten zu unternehmen. Freilich war die Beobachtung
dieses lichtschwachen Weltkörpers mit mancherlei Schwierigkeiten verknüpft, und die von \textsc{Olbers},
\textsc{Oriani}, \textsc{David} und auf der hiesigen Sternwarte angestellten Beobachtungen sind die einzigen, die
im ganzen verflossenen Jahre gemacht wurden. Sehr erwünscht war ihm daher die zu \emph{Brera} am 30.
August 1804 beobachtete Opposition der Pallas, zu welcher Zeit der mittlere Fehler der VII. Elemente
in der Länge $-7'\,28''3$ und in der Breite $+2'\,14''7$ betrug, und er gründete hierauf und auf alle vorhandene
frühere Beobachtungen folgende VIII. Elemente der Pallas:
\SourceLeaf{268}
\begin{center}
\begingroup\fontsize{9}{11}\selectfont\setlength{\tabcolsep}{3pt}\renewcommand{\arraystretch}{1.1}
\begin{tabular}{lr}
Epoche Seeberger Meridian  & $\left\{\begin{array}{lr}1803\ .&221^\circ\ 31'\ 23''2\\1804\ .&299\ 58\ 38.1\end{array}\right.$ \\
tägliche tropische Bewegung . . . & $771''6802$ \\
Sonnenferne 1803 . . . \raisebox{.3ex}{.} . . . & $301^\circ\ 1'\ 44''1$ \\
aufsteigender Knoten 1803 . . . . & $172^\circ\ 29'\ 6''8$ \\
Excentricität . . . . . . . . . & $0.246101$ \\
Logarithmus der halben Axe . . . & $0.4417647$ \\
Neigung der Bahn . . . . . . . & $34^\circ\ 37'\ 43''2$ \\
\end{tabular}\endgroup
\end{center}

\indent Hiernach berechnete Dr. \textsc{Gauss} ferner folgende Ephemeride für den geocentrischen Lauf der \emph{Pallas} [von 1805 Juli 28 bis 1806 April 30 und gibt die Örter in entsprechender Weise an wie die der Ceres für dieselben Jahre, d. B. S. 262.]

—\quad—\quad—

Dr. \textsc{Gauss} schrieb uns, dass sich diese Ephemeride noch einmal auf rein elliptische Elemente gründe, indem er glaube, dass es bei der Pallas jetzt noch nicht zweckmässig sei, auf die Störungen Rücksicht zu nehmen, da theils alle bisherige Beobachtungen sich noch ganz gut durch eine reine Ellipse darstellen lassen, theils die Bahn der Pallas, bei der kurzen Dauer ihrer Erscheinung, noch keine so genaue Bestimmung verstattete, um hoffen zu dürfen, durch die Entwickelung aller Perturbations-Gleichungen eine Genauigkeit zu erhalten, die für jene mühevolle Arbeit belohnen könnte. —\quad—\quad—\quad—\quad—\quad— Wäre die Masse der Ceres etwas beträchtlicher, so könnte diese, theils wegen der geringen gegenseitigen Entfernung, theils wegen der beinahe gleichen mittlern Bewegung beider Planeten, bedeutende Störungen in der Pallas-Bahn bewirken.

—\quad—\quad—

Inspector \textsc{Harding}, dem Dr. \textsc{Gauss} sogleich obige neue Ephemeride der Pallas mittheilte, beschäftigte sich mit der so verdienstlichen Arbeit, eine Sternkarte für ihren Lauf zu entwerfen. Leider wird es ihm hier oft an Sternbestimmungen fehlen, da er in diesen südlichen Zonen nicht so von der Histoire céleste unterstützt wird, als es von dieser schätzbaren Sammlung ausserdem der Fall ist.

\hfill[\textsc{von Zach.}]
